% Section 11: Integration with the Festival Platform

\section{Integration with the Festival Platform}
\label{sec:platform_integration}

The Vendor Engine is designed to interface with the larger festival management ecosystem, maintaining compatibility with the central Django monolith, PostgreSQL database, and WebSocket infrastructure. 

\subsection{Platform Integration Interfaces}
The application exposes three primary integration interfaces, combining configuration ingestion, real-time message simulation, and batch synchronizations.

\subsubsection{Static Configuration Ingestion (\code{stalls.json})}
At application startup, the engine reads stall mappings and menu definitions from a static configurations file. This file represents the list of stalls registered in the backend PostgreSQL database:
\begin{lstlisting}
{
  "stalls": [
    { "stallId": "S01", "stallName": "Chaat Corner" },
    { "stallId": "S02", "stallName": "Momo Point" }
  ]
}
\end{lstlisting}
The \code{MockDataLoader} parses this payload to initialize the memory map of stalls. If this file is missing, the application falls back to a built-in default list to allow execution, logging the error to \code{log/app.log}.

\subsubsection{Simulated Real-Time Ingest Stream}
To simulate the WebSocket orders generated by festival-goers ordering via their mobile apps, order payloads are streamed into the bounded ingestion queue. The JSON payload format matches the Django backend's serialization structure:
\begin{lstlisting}
{
  "orderId": "b1a72d3f-5612-4011-8930-cf2201f3e908",
  "stallId": "S01",
  "priority": "VIP",
  "createdAt": 1731234567890,
  "items": [
    { "itemName": "Pani Puri", "quantity": 2, "unitPrice": 60.0 }
  ]
}
\end{lstlisting}
The local consumers parse these incoming payloads to execute localized queue calculations.

\subsubsection{Reconnection Synchronization Payload (\code{sync\_payload.json})}
When the \code{NetworkMonitorDaemon} transitions from offline back to online, it must synchronize served orders. The application compiles served orders into a batch format and exports it to \code{sync\_payload.json}:
\begin{lstlisting}
{
  "syncedAt": 1731234999999,
  "orders": [
    { "orderId": "b1a72d3f-...", "stallId": "S01", "status": "SERVED", "total": 120.0 }
  ]
}
\end{lstlisting}
In a production deployment, the \code{SyncClient} class would upload this payload via an HTTP \code{POST} request to the Django endpoint \code{/api/stalls/sync/}. Writing this file locally acts as the integration gateway.

\subsection{Local Persistence Architecture (\code{offline\_stalls.ser})}
To achieve zero data loss during network connection drops, the system uses binary Java serialization to snapshot the entire application state. The \code{offline\_stalls.ser} binary file stores the serialized \code{AppState} object graph. This local cache holds stalls, revenue totals, and queues, shielding operations from connection drops.

\subsection{JDBC-Ready DAO Design}
To align with future enhancements, the database operations are structured to follow the Data Access Object (DAO) pattern. 
The \code{StallDataSource} class encapsulates access to the memory-resident stall data. It acts as an abstract data provider, separating query routines from direct mutator methods. 
This interface design is fully ``JDBC-ready'': the memory-backed maps inside \code{StallDataSource} can be swapped with real JDBC connections, queries, and prepared statements without requiring any structural changes to the view or controller classes.

% Source: README.md:17-29, 154-160
% Source: Vendor_Engine_Architecture.md:32-35, 624-625, 647-688
% Source: src/main/java/com/festival/vendorengine/io/NetworkMonitorDaemon.java:208-228
